% Generated from locale/tr/content/counterfactuals/introduction/paradoxes-material.tex; do not hand-edit.


\olfileid[tr]{cnt}{int}{par}

\olsection{Maddi Koşulun Paradoksları}

İngilizcedeki ``if \dots{} then \dots'' ifadelerini simgeleştirmenin bir
yolu olarak $!A\lif !B$'nin kullanılmasını ilk eleştirenlerden biri
C.~I. Lewis'ti. Lewis, Whitehead ile Russell'ın \emph{Principia
Mathematica}'da $\lif$'i ``implies'' diye okudukları maddi koşul kullanımını
eleştiriyordu. Haklı olarak, $\lif$ ``gerektirir'' anlamına gelseydi her
yanlış önerme~$p$'nin, $p$'nin~$q$'yu gerektirdiğini gerektireceğini söyledi;
çünkü $p$ yanlışsa $p\lif(p\lif q)$ doğrudur. Aynı şekilde her doğru
önerme~$q$, $p$'nin~$q$'yu gerektirdiğini gerektirirdi; çünkü $q$ doğruysa
$q\lif(p\lif q)$ doğrudur.

Mantıkçılar elbette \emph{gerektirmenin}, yani mantıksal gerektirmenin bir
bağlaç değil, !!{formula}s ya da ifadeler arasındaki bir bağıntı olduğunu
bilir. Dolayısıyla karışıklığı önlemek için $\lif$'i ``gerektirir'' diye
okumamalıyız.\footnote{Filozoflar Lewis'in belirttiği karışıklıklardan
  kaçınmak için bunu ``yalnızca eğer'' diye okumaya çalışsa da ``$\lif$''i
  ``gerektirir'' diye okumak matematikçiler ve bilgisayar bilimciler arasında
  hâlâ yaygındır.} Böyle okumadığımız sürece Lewis'in özgül kaygısı ortaya
çıkmaz: $p$'nin yanlış bir İngilizce tümceyi temsil ettiğini düşünsek bile
$p$, $q$'yu ``gerektirmez.'' $p\Entails q$ olup olmadığını belirlemek için
\emph{bütün} !!{valuation}s{}i incelemeliyiz; $p$'yi gerçekte yanlış olan bir
tümceyi simgelemek için kullansak bile $p\Entails/ q$ olabilir.

Yine de Lewis'in şu örneği gibi ``eğer \dots{} ise \dots'' tümcelerinde
garip bir şey vardır:
\begin{quote}
Ay yeşil peynirden yapılmışsa $2+2=4$'tür.
\end{quote}
Aşağıdaki çıkarımlarda da benzer bir tuhaflık vardır:
\begin{quote}
  Ay yeşil peynirden yapılmamıştır. Öyleyse ay yeşil peynirden yapılmışsa
  $2+2=4$'tür.

  $2+2=4$'tür. Öyleyse ay yeşil peynirden yapılmışsa $2+2=4$'tür.
\end{quote}
Oysa ``eğer \dots{} ise \dots'' yalnızca $\lif$ olsaydı bu tümce açıkça
doğru, çıkarımlar da açıkça geçerli olurdu.

Başka bir örnek $(!A\lif !B)\lor(!B\lif !A)$ totolojisidir. Bu, gazeteden
gelişigüzel iki bildirme tümcesi~$S$ ve $T$ seçildiğinde ``Eğer $S$ ise $T$,
ya da eğer $T$ ise~$S$'' tümcesinin doğru olması gerektiğini düşündürür.

